\subsection{补算：为什么环平均产生Bessel函数}
以一个垂直Fourier模$\varphi(\vct x)=\varphi_{\vct k}e^{i\vct k_\perp\cdot\vct x}$为例。粒子位置$\vct x=\vct R+\vct r_L(\vartheta)$，固定导心$\vct R$和回旋半径$r_L$，对回旋角平均：
\begin{align}
 \langle\varphi\rangle_\vartheta
 &={1\over2\pi}\int_0^{2\pi}
 \varphi_{\vct k}e^{i\vct k_\perp\cdot\vct R}
 e^{ik_\perp r_L\cos\vartheta}\dd\vartheta\\
 &=J_0(k_\perp r_L)\varphi_{\vct k}e^{i\vct k_\perp\cdot\vct R}.
\end{align}
因此环平均在Fourier空间成为乘以$J_0$，在真实空间则是非局地采样。离子看到的是整个回旋圆上的平均电势，而不是导心处单一点的电势。

对含垂直速度分布的响应再作相应积分，会出现
$\Gamma_0(b)=I_0(b)e^{-b}$。按本节$r_{Li}=\sqrt{T_i/m_i}/\Omega_{ci}$的约定，$b=k_\perp^2r_{Li}^2$。不能把固定回旋半径的$J_0$与Maxwell加权后的$\Gamma_0$当成同一个函数。
小$b$展开为
\begin{equation}
 I_0(b)=1+{b^2\over4}+\cdots,\qquad
 e^{-b}=1-b+{b^2\over2}+\cdots,\qquad
 \Gamma_0=1-b+{3\over4}b^2+\cdots.
\end{equation}
常见Padé近似$(1+b)^{-1}$只在低阶匹配该响应，并不是对任意波数都精确。

\subsection{极化关系的冷离子极限不能直接把$T_i$代零}
本章约化模型中的极化响应包含
$ (\Gamma_0-1)/\tau_i$，其中$\tau_i=T_i/T_e$、$b=\tau_i k_\perp^2\rho_{\rm sound}^2$。当$\tau_i\to0$时，分子与分母同时趋零，必须先展开：
\begin{equation}
 {\Gamma_0-1\over\tau_i}
 ={ -\tau_i k_\perp^2\rho_{\rm sound}^2+O(\tau_i^2)\over\tau_i}
 \longrightarrow-k_\perp^2\rho_{\rm sound}^2.
\end{equation}
回到真实空间，这就是$\rho_{\rm sound}^2\nabla_\perp^2$。所以冷离子极限仍保留极化响应；若先把$\Gamma_0$置为1，再把整个比值删去，就会错误地丢掉涡量--密度耦合。

\paragraph{手算检查。}
对$\phi=\phi_0\cos kx$，有$\nabla_\perp^2\phi=-k^2\phi$。在前面指定的零被动密度模条件下，$N=\rho_{\rm sound}^2\nabla^2\Phi$与势扰动符号相反。这个符号是既定极化与电势约定的结果，不能只凭“正电势应该吸引或排斥什么粒子”替代约化算子的检查。
